\chapter{Hives and Angioedema (Urticaria)}

\section*{Definition}
Hives (urticaria) are raised, itchy, red or skin-colored welts that appear suddenly on the skin. Angioedema is deeper swelling of the subcutaneous and submucosal tissues, often affecting the face, lips, tongue, throat, and genitals. They can occur separately or together.

\section*{What Happens in Your Body}
Both conditions result from mast cell degranulation and release of histamine and other inflammatory mediators. In urticaria, widening of blood vessels and increased vascular permeability in the superficial dermis causes the characteristic wheals. In angioedema, the same process occurs in the deep dermis and subcutaneous tissues. Triggers can be allergic (IgE-mediated) or non-allergic (direct mast cell activation).

\section*{Causes}
\begin{itemize}
\item Allergic reactions: foods (nuts, shellfish, eggs, milk), medications (antibiotics, NSAIDs, aspirin), insect stings, latex
\item Physical triggers: pressure, cold, heat, sunlight, exercise, water, vibration
\item Infections: viral, bacterial, parasitic
\item Autoimmune conditions: thyroid disease, lupus
\item Complement-mediated: hereditary angioedema (C1 inhibitor deficiency)
\item Idiopathic: no identifiable cause
\item Chronic spontaneous urticaria (lasting more than 6 weeks)
\end{itemize}

\section*{When to See a Doctor}
\begin{itemize}
\item Raised, red, itchy welts on the skin
\item Welts that change shape, size, and location within hours
\item Swelling of the face, lips, tongue, throat, hands, feet, or genitals
\item Stinging or burning sensation
\item Individual welts lasting less than 24 hours
\item Angioedema: pain rather than itching
\item Difficulty breathing, throat tightness, or hoarseness (anaphylaxis warning)
\end{itemize}

\section*{Related Symptoms}
\begin{itemize}
\item Anaphylaxis (severe, life-threatening allergic reaction)
\item Allergic rhinitis
\item Asthma
\item Atopic dermatitis (eczema)
\item Contact dermatitis
\end{itemize}

\section*{Self-Care / Home Management}
\begin{itemize}
\item Identify and avoid known triggers
\item Take OTC antihistamines (cetirizine, loratadine, diphenhydramine)
\item Apply cool compresses to reduce itching
\item Wear loose, comfortable clothing
\item Use fragrance-free moisturizers
\item Avoid hot showers, alcohol, and NSAIDs which can worsen symptoms
\item Keep a symptom diary to identify triggers
\end{itemize}

\section*{How Your Doctor Will Evaluate You}
\begin{itemize}
\item Clinical history and physical examination
\item Skin prick testing or specific IgE tests for suspected allergens
\item Physical challenge tests for physical urticarias
\item Complete blood count, ESR, CRP
\item Thyroid antibodies and thyroid function tests
\item Complement levels (if hereditary angioedema suspected)
\item Skin biopsy for chronic or atypical cases
\end{itemize}

\section*{Treatment Options}
\begin{itemize}
\item Second-generation antihistamines as first-line therapy
\item Dose escalation of antihistamines up to 4x standard dose
\item Add H2 blocker (ranitidine, famotidine)
\item Leukotriene receptor antagonists (montelukast)
\item Cyclosporine for refractory chronic urticaria
\item Omalizumab (anti-IgE) for chronic spontaneous urticaria
\item Epinephrine auto-injector for angioedema with anaphylaxis
\item Corticosteroids for acute severe flare-ups
\end{itemize}

\section*{References}
\begin{thebibliography}{99}
\bibitem{ref1} Zuberbier T, Aberer W, Asero R, et al. "The EAACI/GA2LEN/EDF/WAO Guideline for the Definition, Classification, Diagnosis and Management of Urticaria." Allergy. 2018;73(7):1393--1414.
\bibitem{ref2} Kaplan AP, Greaves MW. "Angioedema." J Am Acad Dermatol. 2005;53(3):373--388.
\bibitem{ref3} Grattan CEH, Sabroe RA, Greaves MW. "Chronic Urticaria." J Am Acad Dermatol. 2002;46(5):645--657.
\bibitem{ref4} Kaplan AP. "Chronic Urticaria and Angioedema." N Engl J Med. 2002;346(3):175--179.
\bibitem{ref5} Maurer M, Weller K, Bindslev-Jensen C, et al. "Unmet Clinical Needs in Chronic Spontaneous Urticaria." J Eur Acad Dermatol Venereol. 2011;25(6):637--643.
\bibitem{ref6} Bernstein JA, Lang DM, Khan DA, et al. "The Diagnosis and Management of Acute and Chronic Urticaria: 2014 Update." J Allergy Clin Immunol. 2014;133(5):1270--1277.
\end{thebibliography}
\clearpage
